% Generated draft: requires mathematical and source review.
\documentclass[11pt]{article}
\usepackage[margin=25mm]{geometry}
\usepackage{fontspec}
\setmainfont{Latin Modern Roman}
\usepackage{amsmath,amssymb,mathtools,unicode-math}
\setmathfont{Latin Modern Math}
\usepackage{enumitem,needspace,xurl,hyperref}
\setlength{\parindent}{0pt}
\setlength{\parskip}{5pt}
\setlength{\emergencystretch}{3em}
\newcommand{\N}{\mathbb{N}}
\newcommand{\Z}{\mathbb{Z}}
\newcommand{\Q}{\mathbb{Q}}
\newcommand{\R}{\mathbb{R}}
\newcommand{\C}{\mathbb{C}}
\newcommand{\F}{\mathbb{F}}
\newcommand{\A}{\mathbb{A}}
\newcommand{\PP}{\mathbb{P}}
\newcommand{\E}{\mathbb{E}}
\newcommand{\eps}{\varepsilon}
\newcommand{\dd}{\,\mathrm{d}}
\begin{document}
\section*{UnsolvedMath ID: 30002103}
\textbf{Source problem code:} catalogue entry
\subsection{Definitions and mathematical statement}
Let \(X\) be a singular complex projective variety, and write \(H^n(X,\Z)\) and \(H^n(X,\Q)\) for its singular cohomology groups. The source asks about a filtration denoted by \(L^\bullet H^n(X,\Z)\), called the Leray filtration. A Leray filtration requires an underlying morphism or spectral-sequence construction and an indexing convention; none is specified in the supplied record. Thus \(L^\bullet\) is retained as unresolved source notation.

A rational mixed Hodge structure consists of a finite-dimensional rational vector space with an increasing weight filtration and a decreasing Hodge filtration on its complexification, inducing a pure Hodge structure of weight \(k\) on each weight-graded piece of index \(k\). A sub-mixed-Hodge structure is a subgroup or subspace with the induced filtrations for which inclusion is a morphism of mixed Hodge structures, subject to the chosen integral or rational convention. Write \(\operatorname{Hom}_{\mathrm{MHS}}\) for morphisms of mixed Hodge structures, and \(\Z(-p)\) for the rank-one integral Tate structure of weight \(2p\) and Hodge type \((p,p)\). The treatment of torsion in integral cohomology must be specified to interpret the source's integral formulation.

\textbf{Source question, with unresolved notation retained.} For every singular complex projective variety \(X\), is the Leray filtration on \(H^n(X,\Z)\) a filtration by sub-mixed-Hodge structures? Furthermore, does the proposed \(p\)th Chern-class map have image satisfying
\[
\operatorname{im}\!\left(
K_0(X)\otimes_{\Z}\Q
\xrightarrow{\ c_p\ }
H^{2p}(X,\Q)
\right)
\stackrel{?}{=}
\operatorname{Hom}_{\mathrm{MHS}}\!\left(
\Z(-p),L^pH^{2p}(X,\Z)
\right)\otimes_{\Z}\Q?
\]
The right-hand side is intended to be viewed inside \(H^{2p}(X,\Q)\) by evaluating morphisms on the underlying Tate generator and then using inclusion and rationalization. This interpretation presupposes that the indicated Leray step has the required integral mixed Hodge structure. The source describes the proposed equality as saying that the Chern-class images are precisely the rational Hodge classes in the \(p\)th Leray step. Here rational Hodge classes in a rational mixed Hodge structure \(V\) mean the elements represented by morphisms \(\Q(-p)\to V\).

\textbf{Qualifications required for precision.} The record does not define \(K_0(X)\): on a singular variety, the Grothendieck group of algebraic vector bundles must be distinguished from the Grothendieck group of coherent sheaves. Moreover, the ordinary \(p\)th Chern class is generally not additive for \(p>1\), so the displayed arrow from a rationalized Grothendieck group requires an explicit definition. The record does not specify whether its intended image is the image of a Chern-class operation extended by a stated convention, a rational span of Chern classes, or an additive characteristic-class construction. No replacement by a Chern character is made here. The ranges of \(n\) and \(p\) are not explicitly stated; nonnegative integral indices are suggested by the notation but require confirmation. These omissions prevent a uniquely specified mathematical formulation of the displayed question.
\subsection{Short English statement}
The source asks whether a Leray filtration on the cohomology of every singular complex projective variety respects mixed Hodge structures, and whether Chern classes give exactly the rational Hodge classes in the indicated filtration step. The construction of the filtration, the meaning of \(K_0(X)\), and the proposed Chern-class map require clarification. This is a statement-only draft requiring human mathematical review.
\subsection{Sources}
\begin{itemize}
\item Supplied local UnsolvedMath registry record, permanent ID \(30002103\), source problem code \texttt{catalogue entry}, in the collection \emph{Oberwolfach Reports Open Problems}. No individual proposer is supplied.
\item Supplied report citation: \emph{Classical Algebraic Geometry (2013)}, Oberwolfach Reports, pp.~1845--1893, DOI \texttt{10.4171/owr/2012/30}. \url{https://doi.org/10.4171/owr/2012/30}. The record describes Oberwolfach Reports as documenting mathematical workshops and their open-problem discussions. Its proposed year and year tag are 2012, whereas the supplied citation title gives 2013; this discrepancy is unresolved.
\item Supplied literature source: \emph{Hodge conjecture via singular varieties}. \url{https://arxiv.org/abs/2509.25273}.
\item Supplied literature source labeled ``original OWR report'': \url{https://doi.org/10.4171/owr/2012/30}.
\item Literature assessment dated exactly \texttt{2026-08-21}: the supplied curation labels the problem open and states that the singular cohomological Hodge conjecture has partial results and active modern work, while the stated Chern-class characterization is not proved in general, including a broad singular threefold setting mentioned in that assessment. This is supplied dated curation derived from the OWR record, not an independently verified literature assessment. The source explicitly calls for independent verification of time-sensitive status and cited claims.
\item The registry labels the statement \texttt{corrected\_verified} and reports comparison with the official Oberwolfach report, with notation, context, and question form normalized without changing the source-level claim. This is a supplied verification assertion; neither the report nor the literature source has been independently inspected for this draft.
\end{itemize}
\end{document}
